\chapter{Diskussion}

Social engineering som ett cyberhot har bevisats utmana uppdelningen mellan teknisk och mänsklig säkerhet.
Medan traditionella cybersäkerhetsmodeller ofta fokuserar på systemarkitektur, nätverksskydd och tekniska sårbarheter,
utnyttjar social engineering i första hand organisatoriska strukturer och mänskliga beslutsprocesser.
Detta innebär att angreppsytan för social engineering-relaterade attacker inte endast utgörs av tekniska system som servrar,
protokoll eller autentiseringssystem, utan även av arbetsrutiner, hierarkier och sociala relationer. I motsats till rent tekniska attacker, där sårbarheter ofta kan identifieras och åtgärdas,
är den mänskliga faktorn i social engineering-attacker mer dynamisk och svår att standardisera. En organisation
kan implementera flerfaktorautentisering (M2FA), SOAR-system och AI-baserad detektion, men om en anställd ändå blir manipulerad av en hotaktör
att lämna ut känsliga data eller utföra en skadlig handling, förlorar systemen sin effektivitet. För organisationer betyder detta
att säkerhetsnivån är beroende av både tekniska kontroller och mänskligt omdöme, där den svagaste länken kan avgöra utfallet.

Utvecklingen av AI har endast förstärkt problemet. Medan AI möjliggör mer avancerad detektion och analys, har samma teknologi också börjat användas av hotaktörer
för att skapa mer avancerade attacker, som djupfejk-baserade bedrägerier. Resultatet leder till en teknologisk symmetri där försvar och angrepp utvecklas pararellt. 
Här uppstår det frågan om avancerad teknik kan lösa problemet, eller om den endast förändrar dess form. Organisatorisk säkerhetskultur kopplar sig också till problemet.
I teorin kan säkerhetsutbildning och medvetandeprogram minska risken för lyckad manipulation. I praktiken orsakar överexponering för varningar och säkerhetsrutiner säkerhetsutmattning,
där anställda rutinmässigt godkänner förfrågningar utan noggrann granskning. Detta illusterar att även välmenade säkerhetsåtgärder kan orsaka motsatt effekt om de inte balanseras med 
användarbarhet och organisatoriska rutiner.

Social engineering tar fram också en bredare fråga om ansvarsfördelning. Om en attack lyckas genom manipulation av en individ, ligger ansvaret på den enskilda användaren,
organisationens rutiner eller systemdesignen? Oftast kritiseras den ``bristfälliga användaren``, men forskning har visat att säkerhetsdiskurser ibland individualiserar problemet,
medan sårbarheten ofta är strukturell. Därför måste det tas hänsyn till att social engineering-relaterade incidenter är ett symptom på organisatoriska och kommunikativa brister,
snarare än enbart individuella misstag. Fallstudierna visar därutöver att konsekvenserna varierar i stor grad beroende på kontext. I ekonomiskt motiverade attacker är målet direkt
finansiell vinning, medan politiskt motiverade attacker syftar till att påverka samhälleligt förtroende och informationsmiljöer. Variationen bakom motivet av attacken indikerar att
social engineering inte är en homogen angreppsform utan en flexibel metod som anpassas efter målsättning.